@@M0@@

This relation itself shows the additivity of the map @@M1@@, which can therefore be extended to a homomorphism from @@M2@@ to @@M3@@:

\sgaStableTarget{sga6.e06.rem.1-13}\textbf{Remark 1.13.} We can take as a basis of @@M4@@ over @@M5@@ any sequence of powers of @@M6@@; @@M7@@, since @@M8@@ is invertible, and such a sequence therefore consists of elements linearly independent over @@M9@@; moreover, the other powers of @@M10@@ can be expressed explicitly as linear combinations of these with coefficients in @@M11@@, by equality (1.9.1) and the invertibility of @@M12@@.

Put @@M13@@ (@@M14@@ is interpreted as the ``class of a hyperplane section''). We can then take as a basis of @@M15@@ over @@M16@@ the elements @@M17@@. To obtain the equation of dependence of @@M18@@ we consider the exact sequence:
@@M19@@
and the exact sequence obtained by passing to the duals:
@@M20@@
We then observe that @@M21@@, from which we deduce:
@@M22@@
that is:
\sgaStableTarget{sga6.e06.eq.1-13-1}@@M23@@

We obtain an analogous result by introducing @@M24@@.

Finally, note that if @@M25@@, and hence @@M26@@, we have @@M27@@, and equality (1.13.1) reduces to:
